\begin{theorem}[Positive rescaling of Kronecker factors]%
    \label{thm:kronecker_factor_positive_rescaling}
    \lean{Matrix.exists_smul_posSemidef_of_kronecker_posSemidef}
    \lean{Matrix.eq_zero_of_forall_star_dotProduct_mulVec_eq_zero}
    \lean{Matrix.posSemidef_of_forall_star_dotProduct_mulVec_nonneg}
    \lean{Matrix.star_dotProduct_kronecker_mulVec_prod}
    \lean{Matrix.smul_kronecker_smul_inv_eq}
    \leanok
    Let $A$ and $B$ be nonzero complex square matrices whose Kronecker
    product $A\otimes B$ is positive semidefinite. Then there is a scalar
    $c\ne0$ with $cA\ge0$ and $c^{-1}B\ge0$. Since
    $(cA)\otimes(c^{-1}B)=A\otimes B$, the rescaled factors represent the
    same product.
\end{theorem}

\begin{proof}
    \leanok
    On a product vector $x\otimes y$ the quadratic form of $A\otimes B$
    factors, so
    $(x^\dagger A\,x)(y^\dagger B\,y)\ge0$ for all
    vectors $x$ and $y$.
    By polarization a complex matrix is determined by its quadratic form;
    since $A$ and $B$ are nonzero there are vectors with
    $a=v_0^\dagger A\,v_0\ne0$ and $c=w_0^\dagger B\,w_0\ne0$. Every value
    $x^\dagger(cA)\,x=(x^\dagger A\,x)\,c$ is non-negative, and over the
    complex numbers a non-negative quadratic form already forces the matrix
    to be Hermitian, hence positive semidefinite. For the second factor,
    $ac>0$ and $a\,(y^\dagger B\,y)\ge0$ give
    $c^{-1}(y^\dagger B\,y)=(a\,(y^\dagger B\,y))(ac)^{-1}\ge0$.
\end{proof}

\begin{lemma}[A matrix and its negative cannot both be positive]%
    \label{lem:possemidef_and_neg_possemidef}
    \lean{Matrix.PosSemidef.eq_zero_of_neg_posSemidef}
    \leanok
    If $X$ and $-X$ are positive semidefinite complex matrices, then $X=0$.
\end{lemma}

\begin{proof}
    \leanok
    The quadratic form of $X$ is both non-negative and nonpositive, hence
    vanishes identically. Polarization gives $X=0$.
\end{proof}

\begin{definition}[Finite Kronecker product]%
    \label{def:finite_kronecker_product}
    \lean{Matrix.finKronecker}
    \leanok
    For square matrices $A_n$ indexed by $n=0,\ldots,N-1$, possibly with
    different index types, their finite Kronecker product is the matrix with
    entries
    \begin{align}
        \left(\bigotimes_{n=0}^{N-1}A_n\right)_{x,y}
        &=\prod_{n=0}^{N-1}(A_n)_{x_n,y_n}.
        \notag
    \end{align}
\end{definition}

\begin{lemma}[Elementary finite Kronecker identities]%
    \label{lem:finite_kronecker_elementary}
    \lean{Matrix.finKronecker_ne_zero}
    \lean{Matrix.finKronecker_update_smul}
    \leanok
    \uses{def:finite_kronecker_product}
    A finite Kronecker product of nonzero matrices is nonzero. Replacing one
    factor $A_i$ by $cA_i$ multiplies the whole finite product by $c$.
\end{lemma}

\begin{proof}
    \leanok
    Choose a nonzero entry in each factor and evaluate the product at the
    resulting row and column. The rescaling identity follows immediately from
    the entry formula in Definition~\ref{def:finite_kronecker_product}.
\end{proof}

\begin{theorem}[Uniqueness of the positive rescaling phase]%
    \label{thm:positive_rescaling_phase_unique}
    \lean{Complex.ne_zero_of_norm_eq_one}
    \lean{Complex.unitsPhase}
    \lean{Complex.phase}
    \lean{Complex.phase_smul_posSemidef}
    \lean{Circle.eq_of_smul_posSemidef}
    \lean{Matrix.exists_pos_real_smul_eq_of_smul_posSemidef}
    \leanok
    \uses{lem:possemidef_and_neg_possemidef}
    Let $M\ne0$. If $c_1,c_2\ne0$ and both $c_1M$ and $c_2M$ are positive
    semidefinite, then $c_1=t c_2$ for some positive real number $t$.
    Every complex scalar of unit modulus is nonzero. In particular, every
    nonzero scalar $c$ has a unit-modulus part $u=c/|c|$, positivity of $cM$
    implies positivity of $uM$, and there is at most one such unit scalar when
    $M\ne0$.
\end{theorem}

\begin{proof}
    \leanok
    \uses{lem:possemidef_and_neg_possemidef}
    Hermiticity of $c_1M$ and $c_2M$ gives
    $c_2\overline{c_1}=c_1\overline{c_2}$, so $c_1/c_2$ is real. It cannot be
    negative: otherwise $c_1M$ and its negative would both be positive
    semidefinite, forcing $c_1M=0$. Thus the quotient is positive.
    Multiplication by $|c|^{-1}>0$ replaces $c$ by its unit-modulus part
    without changing positivity. Two admissible unit scalars differ by a
    positive real number of modulus one, and are therefore equal.
\end{proof}

\begin{theorem}[Positive rescaling of a finite Kronecker product]%
    \label{thm:finite_kronecker_positive_rescaling}
    \lean{Matrix.exists_pi_smul_posSemidef_of_finKronecker_posSemidef}
    \leanok
    \uses{def:finite_kronecker_product,
        lem:finite_kronecker_elementary,
        thm:kronecker_factor_positive_rescaling}
    Let $N\geq1$, and let $A_0,\ldots,A_{N-1}$ be nonzero complex square
    matrices, with possibly different dimensions. If their finite Kronecker
    product is positive semidefinite, then there are nonzero scalars
    $c_0,\ldots,c_{N-1}$ such that
    \begin{align}
        \prod_{n=0}^{N-1}c_n&=1,
        \notag\\
        c_nA_n&\geq0 \quad (0\leq n<N).
        \notag
    \end{align}
\end{theorem}

\begin{proof}
    \leanok
    \uses{lem:finite_kronecker_elementary,
        thm:kronecker_factor_positive_rescaling}
    The assertion is proved by induction on $N$. Separate the first factor
    from the remaining Kronecker product and apply
    Theorem~\ref{thm:kronecker_factor_positive_rescaling}. Absorb the inverse
    of the first rescaling coefficient into one factor of the tail, then
    apply the induction hypothesis. The resulting coefficients are nonzero,
    their product is one, and every rescaled factor is positive semidefinite.
\end{proof}

\begin{lemma}[One-sided product obstruction in a full matrix algebra]%
    \label{lem:matrix_one_sided_product_obstruction}
    \lean{Matrix.submodule_sup_ne_top_of_mul_eq_zero}
    \leanok
    Let $\mathcal A_1,\mathcal A_2\subseteq\MN{D}$ be nonzero linear
    subspaces. If $XY=0$ for every $X\in\mathcal A_1$ and
    $Y\in\mathcal A_2$, then
    $\mathcal A_1+\mathcal A_2\ne\MN{D}$.
    No vanishing condition on products in the reverse order is required.
    This is the matrix-algebra obstruction used in
    \cite[Appendix~C.2, lines~1465--1470]{Cirac2016MPDO_arXiv}.
\end{lemma}

\begin{proof}
    \leanok
    Choose nonzero matrices $X\in\mathcal A_1$ and
    $Y\in\mathcal A_2$, together with nonzero entries $X_{ip}$ and
    $Y_{qj}$. If the two subspaces spanned the full matrix algebra, write
    the matrix unit $E_{pq}=E_1+E_2$ with
    $E_1\in\mathcal A_1$ and $E_2\in\mathcal A_2$. The one-sided product
    hypothesis gives
    \begin{align}
        XE_{pq}Y&=X(E_1Y)+(XE_2)Y=0.
        \notag
    \end{align}
    Its $(i,j)$ entry is $X_{ip}Y_{qj}\ne0$, a contradiction.
\end{proof}

\begin{definition}[Nonzero sector support and recurrence]%
    \label{def:mpdo_eta_recurrent_support}
    \lean{MPOTensor.IsSectorEdge}
    \lean{MPOTensor.SectorReaches}
    \lean{MPOTensor.IsRecurrentSupport}
    \leanok
    \uses{def:mpdo_cyclic_eta_tensor_product}
    The nonzero support graph of a neighboring-operator family has a directed
    edge $k\to h$ precisely when $\eta_{k,h}\ne0$. Its support is recurrent
    when every such edge belongs to a directed cycle, equivalently when $h$
    can reach $k$ whenever $k\to h$ is an edge.
\end{definition}

\begin{definition}[Sector-coordinate virtual matrix families]%
    \label{def:mpdo_sector_virtual_matrix_families}
    \lean{MPOTensor.PhysicalSectorFactorization.sectorCoordinateMatrixFamily}
    \lean{MPOTensor.PhysicalSectorFactorization.sectorVirtualMatrix}
    \lean{MPOTensor.PhysicalSectorFactorization.sectorVirtualMatrixFamily}
    \leanok
    \uses{def:mpdo_physical_sector_factorization}
    In sector coordinates, let $A_{k;x,y}$ be the virtual matrix obtained
    from the $(x,y)$ entry within sector $k$. The sector-coordinate family
    contains all transformed physical matrices, while the virtual-matrix
    family consists of all matrices $A_{k;x,y}$.
\end{definition}

\begin{theorem}[Sector spanning from injectivity]%
    \label{thm:mpdo_sector_virtual_matrix_spanning}
    \lean{MPOTensor.PhysicalSectorFactorization.changePhysicalBasis_inverse_apply_eq_sum}
    \lean{MPOTensor.PhysicalSectorFactorization.sectorCoordinateMatrixFamily_span_eq_top}
    \lean{MPOTensor.PhysicalSectorFactorization.sectorCoordinateTensor_mem_span_sectorVirtualMatrixFamily}
    \lean{MPOTensor.PhysicalSectorFactorization.sectorVirtualMatrixFamily_span_eq_top}
    \leanok
    \uses{def:mpdo_sector_virtual_matrix_families}
    Let ${\cal K}$ be injective and let
    \begin{align}
        U\kappa_{\beta,\alpha}U^\dagger
        &=\bigoplus_k(l_k)_\beta\otimes(r_k)_\alpha
        \notag
    \end{align}
    be a physical-sector factorization. Then the virtual matrices
    $A_{k;x,y}$ obtained from all within-sector physical entries span the full
    virtual matrix algebra $\MN{D}$.
\end{theorem}

\begin{proof}
    \leanok
    Since $U$ is unitary, every original physical matrix of ${\cal K}$ is a
    linear combination of the matrices in the transformed physical
    coordinates. The off-diagonal sector entries in those coordinates vanish,
    while each diagonal sector entry is one of the matrices $A_{k;x,y}$.
    Injectivity says that the original physical matrices span $\MN{D}$, and
    the conclusion follows.
\end{proof}

\begin{theorem}[Triangle closure from sector spanning]%
    \label{thm:mpdo_sector_support_triangle_closure}
    \lean{MPOTensor.PhysicalSectorFactorization.sectorVirtualMatrix_mul_apply}
    \lean{MPOTensor.PhysicalSectorFactorization.%
      sectorVirtualMatrix_mul_eq_zero_of_neighboringOperator_eq_zero}
    \lean{MPOTensor.PhysicalSectorFactorization.%
      exists_two_edge_return_of_neighboringOperator_ne_zero_of_endpoints}
    \lean{MPOTensor.PhysicalSectorFactorization.exists_two_edge_return_of_neighboringOperator_ne_zero}
    \lean{MPOTensor.PhysicalSectorFactorization.neighboringOperator_reaches_of_ne_zero}
    \leanok
    \uses{def:mpdo_sector_virtual_matrix_families,
        def:mpdo_minimal_neighboring_operator,
        thm:trace_nondeg}
    For a physical-sector factorization, let $A_{k;x,y}\in\MN{D}$ be the
    virtual matrix obtained from the physical matrix entry $(x,y)$ within
    sector $k$. Suppose that all such matrices span $\MN{D}$ and that both
    endpoint sectors of every nonzero neighboring operator contain a nonzero
    matrix of this form. Then every nonzero $\eta_{k,h}$ admits a two-edge
    return: there is a sector $j$ for which
    \begin{align}
        \eta_{h,j}&\ne0,
        \notag\\
        \eta_{j,k}&\ne0.
        \notag
    \end{align}
    In particular, the nonzero neighboring support is recurrent.
\end{theorem}

\begin{proof}
    \leanok
    \uses{thm:trace_nondeg}
    A nonzero entry of $\eta_{k,h}$, together with nonzero sector matrices at
    $k$ and $h$, gives matrices $A$ from sector $k$ and $B$ from sector $h$
    with $AB\ne0$. By nondegeneracy of the trace pairing, choose $Y$ with
    $\tr(ABY)\ne0$. Expanding $Y$ in the spanning family, some sector matrix
    $C$ from a sector $j$ satisfies $\tr(ABC)\ne0$. Cyclicity of the trace
    gives $BC\ne0$ and $CA\ne0$. The sector multiplication identity then
    forces $\eta_{h,j}\ne0$ and $\eta_{j,k}\ne0$.
\end{proof}

\begin{lemma}[Equivalent forms of sector reachability]%
    \label{lem:mpdo_eta_reachability_equivalence}
    \lean{TNLean.Algebra.DirectedWalk.reaches_iff_reflTransGen}
    \lean{MPOTensor.sectorReaches_iff_directedWalkReaches}
    \lean{MPOTensor.isRecurrentSupport_iff_directedWalkReturns}
    \leanok
    A vertex $h$ is reachable from $k$ in the nonzero sector support if and
    only if there is a finite directed walk from $k$ to $h$. Consequently,
    the support is recurrent if and only if every directed edge admits a
    directed return walk.
\end{lemma}

\begin{proof}
    \leanok
    A walk gives an element of the reflexive transitive closure by induction
    on its edges. Conversely, append the final edge at each induction step in
    the reflexive transitive closure.
\end{proof}

\begin{lemma}[Finite cyclic form of a closed directed walk]%
    \label{lem:directed_closed_walk_finite_cycle}
    \lean{TNLean.Algebra.DirectedWalk.closedConsVertices}
    \lean{TNLean.Algebra.DirectedWalk.closedConsVertices_edge}
    \lean{TNLean.Algebra.DirectedWalk.weight_closedCons_eq_fin_prod}
    \leanok
    Every nonempty closed directed walk of length $N$ determines a cyclic
    family of vertices $v_0,\ldots,v_{N-1}$. Each $v_n\to v_{n+1}$ is an
    edge, where $v_N=v_0$, and for any edge weights in a commutative monoid,
    $\operatorname{wt}(w)=\prod_{n=0}^{N-1}\kappa_{v_n,v_{n+1}}$.
\end{lemma}

\begin{proof}
    \leanok
    Enumerate the sources of the successive edges. The last edge ends at the
    initial vertex because the walk is closed. The product identity follows
    by induction on the number of edges.
\end{proof}

\begin{theorem}[Closed-walk criterion for vertex phases]%
    \label{thm:directed_cycle_coboundary}
    \lean{TNLean.Algebra.DirectedWalk.exists_vertex_of_closedWalk_weight_eq_one}
    \leanok
    Let $E$ be a directed relation on a set $V$, and suppose that every edge
    $a\to b$ admits a directed return walk from $b$ to $a$. Let $G$ be a
    group and assign a weight $\kappa_{a,b}\in G$ to each ordered
    pair of vertices. If the product of the edge weights along every closed
    directed walk is one, then there are vertex weights $z_a\in G$ such that
    $\kappa_{a,b}=z_a^{-1}z_b$ whenever $a\to b$ is an edge. This is the
    graph-theoretic coboundary step needed to make the phase choices in
    \cite[Appendix~C.2, lines~1446--1450]{Cirac2016MPDO_arXiv} coherent across
    a recurrent support component.
\end{theorem}

\begin{proof}
    \leanok
    Reachability is an equivalence relation because every edge has a return
    walk. Choose one vertex in each reachability class and, for every vertex
    $a$, choose a walk from the class representative to $a$. Define $z_a$ as
    the product of the edge weights along this walk. If two such walks have
    the same endpoints, append a common return walk; the two resulting closed
    walks have weight one, so the original walk weights agree. Comparing the
    chosen walk to $b$ with the chosen walk to $a$ followed by the edge
    $a\to b$ gives $z_b=z_a\kappa_{a,b}$.
\end{proof}

\begin{definition}[Horizontal bond fiber of a sector cycle]%
    \label{def:mpdo_eta_cycle_fiber}
    \lean{MPOTensor.CycleFiber}
    \lean{MPOTensor.cycleFiberToSectorFiber}
    \leanok
    \uses{def:mpdo_cyclic_eta_tensor_product}
    For a cyclic sector assignment $(k_n)_{n\in\mathbb Z/N\mathbb Z}$, the
    horizontal bond fiber at $n$ is
    \begin{align}
        H_n&=B^R_{k_n}\times B^L_{k_{n+1}}.
        \notag
    \end{align}
    Shifting the left component by one site identifies
    $\prod_n H_n$ with the vertex-indexed sector fiber used in the cyclic
    neighboring-operator product.
\end{definition}

\begin{lemma}[A cyclic block is a finite Kronecker product]%
    \label{lem:mpdo_eta_cycle_fin_kronecker}
    \lean{MPOTensor.cyclicEtaTensorProduct_submatrix_eq_finKronecker}
    \leanok
    \uses{def:finite_kronecker_product,
        def:mpdo_cyclic_eta_tensor_product,
        def:mpdo_eta_cycle_fiber}
    Reindexing a cyclic neighboring-operator block along its horizontal bond
    fiber gives the finite Kronecker product
    $\bigotimes_{n=0}^{N-1}\eta_{k_n,k_{n+1}}$.
\end{lemma}

\begin{proof}
    \leanok
    Expand both sides entrywise. The shifted left component of the bond fiber
    at site $n+1$ is precisely the left index paired with the right component
    at site $n$.
\end{proof}

\begin{theorem}[Positive rescaling around one sector cycle]%
    \label{thm:mpdo_eta_cycle_positive_rescaling}
    \lean{MPOTensor.exists_pi_smul_posSemidef_of_cyclicEtaTensorProduct_posSemidef}
    \leanok
    \uses{def:mpdo_cyclic_eta_tensor_product,
        lem:mpdo_eta_cycle_fin_kronecker,
        thm:finite_kronecker_positive_rescaling}
    Let $N\geq1$, and fix a cyclic sector assignment
    $k_0,\ldots,k_{N-1}$ such that every operator
    $\eta_{k_n,k_{n+1}}$ is nonzero. If the corresponding cyclic tensor
    product is positive semidefinite, then there are nonzero scalars $c_n$
    satisfying
    \begin{align}
        \prod_{n=0}^{N-1}c_n&=1,
        \notag\\
        c_n\eta_{k_n,k_{n+1}}&\geq0 \quad (0\leq n<N).
        \notag
    \end{align}
    This is a choice on one fixed cycle. It does not assert that choices made
    on two different cycles agree on a common edge.
\end{theorem}

\begin{proof}
    \leanok
    \uses{lem:mpdo_eta_cycle_fin_kronecker,
        thm:finite_kronecker_positive_rescaling}
    Reindex the cyclic sector block by the shared horizontal-bond indices.
    Its matrix entries then become the finite Kronecker product of the
    consecutive neighboring operators. Apply
    Theorem~\ref{thm:finite_kronecker_positive_rescaling}.
\end{proof}

\begin{corollary}[Positive rescaling of the concrete sector operators around
    one cycle]%
    \label{cor:mpdo_sector_eta_cycle_positive_rescaling}
    \lean{MPOTensor.exists_pi_smul_posSemidef_of_sectorEta_cycle}
    \leanok
    \uses{def:mpdo, def:injective, def:mpdo_eta_structure,
        def:mpdo_three_site_closure, def:mpdo_sector_eta,
        thm:mpdo_eta_cycle_positive_rescaling,
        thm:mpdo_cyclic_eta_posSemidef}
    Let ${\cal K}$ be an injective MPDO, let $R$ give a three-site closure,
    and fix an $\eta$-structure and a nonzero tail entry. Let $N\geq1$. For a
    cyclic sector assignment $k_0,\ldots,k_{N-1}$, suppose that every
    concrete operator $\eta_{k_n,k_{n+1}}$ obtained from these data is
    nonzero. Then there are nonzero scalars $c_n$ such that
    \begin{align}
        \prod_{n=0}^{N-1}c_n&=1,
        \notag\\
        c_n\eta_{k_n,k_{n+1}}&\geq0 \quad (0\leq n<N).
        \notag
    \end{align}
    This is the specialization of
    Theorem~\ref{thm:mpdo_eta_cycle_positive_rescaling} to the inverse-map
    sector operators; positivity of the cyclic product follows from the MPDO
    property.
\end{corollary}

\begin{proof}
    \leanok
    \uses{thm:mpdo_eta_cycle_positive_rescaling,
        thm:mpdo_cyclic_eta_posSemidef}
    The fixed-sector cyclic product is positive semidefinite by
    Theorem~\ref{thm:mpdo_cyclic_eta_posSemidef}. Apply
    Theorem~\ref{thm:mpdo_eta_cycle_positive_rescaling}.
\end{proof}

\begin{theorem}[The conjugated chain is a congruence]%
    \label{thm:mpdo_conjugated_chain_congruence}
    \lean{MPOTensor.mpo_conjugatePhysical_eq}
    \lean{MPOTensor.mpo_conjugatePhysical_posSemidef}
    \leanok
    \uses{def:mpdo_conjugate_physical}
    For every matrix $U$ on the physical space and every $N\ge1$,
    \begin{align}
        \rho^{(N)}(\widetilde{\cal K})
        &=U^{\otimes N} \rho^{(N)}({\cal K})
          (U^{\otimes N})^\dagger.
        \notag
    \end{align}
    In particular $\rho^{(N)}(\widetilde{\cal K})\ge0$ whenever
    $\rho^{(N)}({\cal K})\ge0$.
\end{theorem}

\begin{proof}
    \leanok
    Expand the closed trace over bond configurations $g$ (with
    $g_{N+1}=g_1$), and each conjugated local entry over its two physical
    indices:
    \begin{align}
        [\rho^{(N)}(\widetilde{\cal K})]_{\sigma,\tau}
        &=\sum_{g}\prod_{n=1}^{N}
            \widetilde\kappa_{g_n,g_{n+1}}^{\,\sigma_n\tau_n}
        \notag\\
        &=\sum_{g}\prod_{n=1}^{N}\sum_{p_n,q_n}
            U_{\sigma_n p_n}
            \kappa_{g_n,g_{n+1}}^{\,p_n q_n}
            \overline{U_{\tau_n q_n}}.
        \notag
    \end{align}
    Distributing the site product over the physical sums factors the two
    $U$-products out of the bond sum, which leaves the matrix entry of
    $U^{\otimes N}\rho^{(N)}({\cal K})(U^{\otimes N})^\dagger$.
    Positivity is preserved by any congruence.
\end{proof}

\begin{theorem}[Positivity of the projected chain blocks]%
    \label{thm:mpdo_cyclic_eta_posSemidef}
    \lean{MPOTensor.cyclicEtaTensorProduct_posSemidef}
    \leanok
    \uses{def:mpdo, def:mpdo_cyclic_eta_tensor_product,
        thm:mpdo_fixed_sector_cyclic_eta,
        thm:mpdo_conjugated_chain_congruence}
    Let ${\cal K}$ be an MPDO whose transformed physical slices satisfy the
    sector factorization hypothesis. Then for every $N\ge1$ and every
    cyclic sector assignment $k_1,\ldots,k_N$ with $k_{N+1}=k_1$,
    \begin{align}
        0
        &\leq
        [Q_{k_1}\otimes\cdots\otimes Q_{k_N}]\,\tilde\sigma
        [Q_{k_1}\otimes\cdots\otimes Q_{k_N}]
        \notag\\
        &=\bigotimes_{n=1}^{N}\eta_{k_n,k_{n+1}}.
        \notag
    \end{align}
    This is the positivity half of the display at
    \cite[Appendix~C.2, lines~1446--1450]{Cirac2016MPDO_arXiv}.
\end{theorem}

\begin{proof}
    \leanok
    \uses{thm:mpdo_fixed_sector_cyclic_eta,
        thm:mpdo_conjugated_chain_congruence}
    The cyclic tensor product $\Omega_k$ of
    Definition~\ref{def:mpdo_cyclic_eta_tensor_product} is a diagonal
    sector block of the basis-conjugated chain by
    Theorem~\ref{thm:mpdo_fixed_sector_cyclic_eta}. The conjugated chain is
    positive semidefinite by
    Theorem~\ref{thm:mpdo_conjugated_chain_congruence} and the MPDO
    property, and a principal submatrix of a positive semidefinite matrix
    is positive semidefinite.
\end{proof}

\begin{theorem}[Positivity of the diagonal and paired neighboring
    operators]%
    \label{thm:mpdo_sector_eta_pair_posSemidef}
    \lean{MPOTensor.sectorEta_self_posSemidef}
    \lean{MPOTensor.sectorEta_kronecker_posSemidef}
    \lean{MPOTensor.exists_smul_posSemidef_sectorEta}
    \lean{MPOTensor.etaOfSectorTensors_self_posSemidef}
    \lean{MPOTensor.etaOfSectorTensors_kronecker_posSemidef}
    \leanok
    \uses{def:mpdo, def:injective, def:mpdo_sector_eta,
        thm:mpdo_cyclic_eta_posSemidef,
        thm:kronecker_factor_positive_rescaling}
    Let ${\cal K}$ be an injective MPDO with the chosen Hayashi
    decomposition and a nonzero tail entry. Then for all sectors $k,h$:
    \begin{enumerate}
        \item $\eta_{k,k}\ge0$;
        \item $\eta_{k,h}\otimes\eta_{h,k}\ge0$;
        \item if $\eta_{k,h}\ne0$ and $\eta_{h,k}\ne0$, there is a scalar
            $c\ne0$ with $c\,\eta_{k,h}\ge0$, $c^{-1}\eta_{h,k}\ge0$, and
            $(c\,\eta_{k,h})\otimes(c^{-1}\eta_{h,k})=\eta_{k,h}\otimes
            \eta_{h,k}$.
    \end{enumerate}
\end{theorem}

\begin{proof}
    \leanok
    \uses{thm:mpdo_cyclic_eta_posSemidef,
        thm:kronecker_factor_positive_rescaling}
    The chain of length one projected to the sector $k$ is $\eta_{k,k}$ up
    to the exchange of its two tensor factors, and the chain of length two
    projected to $(k,h)$ is $\eta_{k,h}\otimes\eta_{h,k}$, both positive by
    Theorem~\ref{thm:mpdo_cyclic_eta_posSemidef}. The third item, including
    the Kronecker-product-preserving equality, is
    Theorem~\ref{thm:kronecker_factor_positive_rescaling}.
\end{proof}

\begin{theorem}[Coherent positive rephasing under recurrent support]%
    \label{thm:mpdo_eta_coherent_positive_choice_recurrent}
    \lean{MPOTensor.exists_vertexPhase_smul_posSemidef}
    \lean{MPOTensor.exists_rephased_inverseMapPhysicalSectorFactorization}
    \lean{MPOTensor.nonempty_etaLocalStructureData_of_recurrentSupport}
    \leanok
    \uses{def:mpdo_physical_sector_rephasing,
        def:mpdo_eta_recurrent_support,
        lem:mpdo_eta_reachability_equivalence,
        lem:directed_closed_walk_finite_cycle,
        thm:positive_rescaling_phase_unique,
        thm:mpdo_eta_cycle_positive_rescaling,
        thm:directed_cycle_coboundary}
    Let $(\eta_{k,h})$ be a neighboring-operator family for which every
    nonempty cyclic tensor product is positive semidefinite. Suppose, in
    addition, that every nonzero edge $k\to h$ admits a directed return path
    from $h$ to $k$. Then there are unit complex numbers $z_k$ such that
    $z_k^{-1}z_h\,\eta_{k,h}\geq0$ for all $k$ and $h$.
    Consequently, if the neighboring operators arise from the inverse-map
    physical-sector factorization of an injective MPDO and its nonzero
    support is recurrent, the sector tensors admit a reciprocal rephasing
    for which all neighboring operators are positive semidefinite. This
    factorization determines the positive commuting nearest-neighbor product
    structure of Proposition~C.8.

    This is a scope-restricted result. Recurrence is not a hypothesis of
    Lemma~C.4 in
    \cite[Appendix~C.2, lines~1406--1450]{Cirac2016MPDO_arXiv}. Its derivation
    for the zero-weight reparameterized inverse-map factorization is treated in
    Theorem~\ref{thm:mpdo_eta_coherent_positive_choice} and recorded in
    \path{docs/paper-gaps/cpgsv17_mpdo_sal_zcl_eta_local_structure.tex}.
\end{theorem}
% The recurrence gap is documented in
% docs/paper-gaps/cpgsv17_mpdo_sal_zcl_eta_local_structure.tex.

\begin{proof}
    \leanok
    \uses{lem:mpdo_eta_reachability_equivalence,
        lem:directed_closed_walk_finite_cycle,
        thm:positive_rescaling_phase_unique,
        thm:mpdo_eta_cycle_positive_rescaling,
        thm:directed_cycle_coboundary,
        def:mpdo_physical_sector_rephasing,
        thm:mpdo_eta_local_structure_of_positive_physical_sector_factorization}
    For a nonzero edge, choose a directed return path and close it with the
    edge. Positive rescaling of the corresponding cyclic tensor product
    supplies a scalar that makes the edge operator positive; retain only its
    unit-modulus part. On any closed directed walk, positive rescaling of the
    associated cyclic tensor product gives phases whose product is one.
    Uniqueness of the unit phase identifies them edge by edge with the chosen
    phases. Hence the chosen edge phases have product one around every closed
    walk. The closed-walk criterion gives
    $\kappa_{k,h}=z_k^{-1}z_h$ on every nonzero edge. Zero edge operators are
    already positive semidefinite. Finally, reciprocal rephasing of the left
    and right sector tensors multiplies each neighboring operator by this
    phase and preserves the physical-slice factorization. The conditional
    eta-local construction then gives the positive commuting product data.
\end{proof}

\begin{theorem}[Coherent positive choice of the neighboring operators]%
    \label{thm:mpdo_eta_coherent_positive_choice}
    \lean{MPOTensor.conjugated_middle_threeSiteClosure}
    \lean{MPOTensor.exists_active_sectorVirtualMatrix_ne_zero}
    \lean{MPOTensor.zeroWeightReparameterizedInverseMapPhysicalSectorFactorization_isRecurrentSupport}
    \lean{MPOTensor.exists_rephase_zeroWeightInverseMap_posSemidef}
    \lean{MPOTensor.exists_positive_inverseMapPhysicalSectorFactorization}
    \lean{MPOTensor.exists_positive_physicalSectorFactorization_of_isSAL}
    \leanok
    \uses{def:mpdo, def:injective, def:is_sal,
        def:mpdo_sector_eta,
        def:mpdo_zero_weight_reparameterized_inverse_map_factorization,
        thm:mpdo_zero_weight_reparameterized_neighboring_operator,
        thm:mpdo_sector_virtual_matrix_spanning,
        thm:mpdo_sector_support_triangle_closure,
        thm:mpdo_eta_coherent_positive_choice_recurrent}
    Let ${\cal K}$ be an injective MPO tensor satisfying SAL. Then it admits
    a physical-sector factorization for which every neighboring operator
    $\eta_{k,h}=\sum_a (r_k)_a\otimes(l_h)_a$ is positive semidefinite.
    Zero-weight Hayashi sectors remain in the physical direct sum, but their
    right tensors may be set to zero without changing the physical slices.
\end{theorem}

\begin{proof}
    \leanok
    \uses{def:mpdo_eta_recurrent_support,
        thm:mpdo_cyclic_eta_posSemidef,
        thm:mpdo_eta_cycle_positive_rescaling,
        def:mpdo_zero_weight_reparameterized_inverse_map_factorization,
        thm:mpdo_zero_weight_reparameterized_neighboring_operator,
        thm:mpdo_sector_virtual_matrix_spanning,
        thm:mpdo_sector_support_triangle_closure,
        thm:positive_rescaling_phase_unique,
        thm:directed_cycle_coboundary,
        thm:mpdo_eta_coherent_positive_choice_recurrent}
    First set the right tensor to zero in every zero-weight sector. The left
    tensor already vanishes there, so this does not alter the factorization
    and removes every support edge incident to such a sector. Every remaining
    endpoint has positive weight. If all virtual matrices in a positive-weight
    sector vanished, then the corresponding middle-site block of the
    three-site closure would vanish; this contradicts its Hayashi expression
    as a nonzero weight times two trace-one sector states. Thus every endpoint
    of a nonzero edge contains a nonzero virtual matrix. Injectivity supplies
    spanning by all sector virtual matrices, and triangle closure gives
    recurrence. The projected chain identity makes every cyclic Kronecker
    product positive semidefinite. Apply
    Theorem~\ref{thm:mpdo_eta_coherent_positive_choice_recurrent} to obtain the
    coherent vertex rephasing.
\end{proof}

\begin{theorem}[Positivity choice with symmetric support]%
    \label{thm:mpdo_eta_positivity_choice_symmetric}
    \lean{MPOTensor.exists_explicitEtaOperators_sectorEta}
    \leanok
    \uses{def:mpdo_explicit_eta_operators,
        thm:mpdo_sector_eta_pair_posSemidef}
    In the setting of Theorem~\ref{thm:mpdo_sector_eta_pair_posSemidef},
    assume in addition that the nonzero neighboring operators occur in
    symmetric pairs: $\eta_{k,h}=0$ forces $\eta_{h,k}=0$. Then there is a
    positive semidefinite family $\eta'_{k,h}$ with $\eta'_{k,k}=\eta_{k,k}$
    and, for every pair of sectors, scalars $c_{k,h}\ne0$ with
    $\eta'_{k,h}=c_{k,h} \eta_{k,h}$ and
    $\eta'_{h,k}=c_{k,h}^{-1} \eta_{h,k}$. In particular every two-site
    block $\eta'_{k,h}\otimes\eta'_{h,k}=\eta_{k,h}\otimes\eta_{h,k}$ of
    the sector-adapted decomposition is unchanged.

    The source asserts at
    \cite[Appendix~C.2, lines~1446--1450]{Cirac2016MPDO_arXiv} that the
    $\eta$'s can be chosen positive semidefinite, without the
    symmetric-support hypothesis. That hypothesis excludes a scalar
    obstruction: a nonzero neighboring operator whose reverse pair
    vanishes is constrained only by cycles through at least three
    sectors, and a coherent choice for the whole family can then fail.
\end{theorem}

\begin{proof}
    \leanok
    \uses{thm:mpdo_sector_eta_pair_posSemidef,
        thm:kronecker_factor_positive_rescaling}
    The diagonal operators are positive by
    Theorem~\ref{thm:mpdo_sector_eta_pair_posSemidef}, and their scalar is
    taken to be one. For a pair of distinct sectors, either both
    neighboring operators vanish and any nonzero scalar works, or both are
    nonzero by symmetric support and
    Theorem~\ref{thm:mpdo_sector_eta_pair_posSemidef} supplies the
    reciprocal pair of scalars. Choosing the scalar once per unordered
    pair makes the two assignments reciprocal, and the block-preserving
    equality of Theorem~\ref{thm:kronecker_factor_positive_rescaling}
    applied to that scalar gives
    $\eta'_{k,h}\otimes\eta'_{h,k}=\eta_{k,h}\otimes\eta_{h,k}$.
\end{proof}
